\cleardoublepage % memastikan bab baru dimulai di halaman ganjil (kanan)
\chapter{PENDAHULUAN}
\label{chap:pendahuluan}
% --- Latar Belakang ---
\section{Latar Belakang}
Ibadah pada bulan puasa mewajibkan umat Muslim untuk menahan lapar dan haus sejak terbit fajar hingga terbenam matahari, sehingga pola makan yang pada hari-hari biasa terbagi menjadi tiga waktu makan utama berubah menjadi hanya dua waktu, yaitu sahur dan berbuka. Perubahan pola makan ini memengaruhi pola harian, termasuk pergeseran jam tidur karena harus bangun pada dini hari untuk beribadah sahur (Mahmudiomo, 2019). Meskipun waktu makan menjadi terbatas, kebutuhan gizi tubuh selama berpuasa tidak berkurang. Pemenuhan gizi seimbang pada waktu sahur dan berbuka menjadi penentu utama agar tubuh tetap bertenaga, tidak lemas, dan mampu menjalankan ibadah serta aktivitas harian secara maksimal (Kementerian Kesehatan RI, 2025).

Pada praktiknya, tidak sedikit orang, termasuk mahasiswa, yang belum mengatur pola makan sahur dan berbukanya dengan memperhatikan keseimbangan gizi. Ketua Program Studi Gizi Universitas Al-Azhar Indonesia mengungkapkan bahwa penurunan performa secara drastis atau \emph{brain fog} yang banyak dialami pekerja maupun mahasiswa selama bulan puasa sebagian besar dipicu oleh kebiasaan sahur yang keliru, misalnya hanya mengandalkan makanan instan tanpa mempertimbangkan kandungan gizinya (Irawan, 2026). Kondisi ini diperparah oleh sempitnya jendela waktu sahur, yang idealnya memberi jeda sekitar dua jam bagi tubuh untuk mencerna makanan sebelum memasuki waktu imsak, sehingga banyak orang cenderung memilih menu seadanya yang praktis namun minim nilai gizi (Kementerian Kesehatan RI, 2025).

Salah satu penyebab mendasar dari pola makan yang tidak seimbang tersebut adalah rendahnya kesadaran dan kebiasaan memantau asupan gizi harian. Banyak orang sulit memahami status gizi dirinya sendiri maupun mengontrol makanan yang dikonsumsi, serta belum terbiasa menghitung atau mencatat kebutuhan kalori dan zat gizi hariannya secara tepat (Anwar, 2021, dalam Jurnal Teknologi Informasi dan Ilmu Komputer Universitas Brawijaya). Di kalangan ahli gizi, praktik pengkajian asupan makan sebenarnya lazim dilakukan melalui metode \emph{24-hour recall}, namun pelaksanaannya secara manual terbukti memakan waktu serta menghasilkan data yang terfragmentasi dan tidak lengkap (Kusumastuty, Handayani, dan Nugroho, 2023). Padahal, selama bulan puasa, kebutuhan untuk memantau asupan gizi menjadi lebih mendesak karena waktu makan yang terbatas membuat setiap pilihan menu sahur maupun berbuka berdampak besar terhadap kondisi tubuh sepanjang hari.

Kebutuhan pemantauan asupan gizi selama bulan puasa tidak dapat dilepaskan dari bagaimana pengguna berinteraksi dengan aplikasi pemantauan itu sendiri. Pada waktu sahur, pengguna umumnya berada dalam kondisi terburu-buru, mengantuk, dan memiliki waktu yang sangat singkat sebelum imsak, sehingga proses pencatatan asupan gizi yang rumit atau memakan banyak langkah berisiko tidak digunakan secara konsisten. Kondisi serupa juga terjadi saat berbuka, ketika pengguna cenderung ingin segera menyantap hidangan setelah seharian menahan lapar dan haus. Oleh karena itu, desain interaksi aplikasi pemantauan gizi selama bulan puasa perlu dirancang secara sederhana dan tidak membebani, dengan mempertimbangkan bagaimana pengguna memperhatikan, memahami, dan memproses informasi dalam kondisi waktu dan energi yang terbatas tersebut (Sharp, Rogers, dan Preece, 2019).

Topik pengembangan aplikasi pemantauan asupan gizi sesungguhnya bukan hal baru dalam penelitian di bidang teknologi kesehatan. Beberapa penelitian sebelumnya telah mengembangkan aplikasi Rekasku, sebuah instrumen pengkajian asupan makan berbasis metode \emph{24-hour recall} yang ditujukan untuk membantu ahli gizi, namun proses pengisiannya masih memerlukan waktu yang relatif panjang (Kusumastuty, Handayani, dan Nugroho, 2023). Penelitian lain mengembangkan aplikasi pemantauan gizi terintegrasi bernama NutriCare yang menggabungkan pemantauan asupan gizi dengan aspek kesehatan mental penggunanya secara umum. Ada pula penelitian yang merancang aplikasi \emph{mobile} untuk \emph{tracking} nutrisi dengan metode \emph{System Development Life Cycle} (SDLC), yang menyoroti keterbatasan aplikasi pemantauan gizi yang ada, yaitu masih sedikit yang berbahasa Indonesia dan sebagian besar hanya berfokus pada penyimpanan data tanpa memperhatikan kemudahan interaksi bagi penggunanya. Di ranah desain antarmuka, penelitian tentang aplikasi Eatplan merancang antarmuka untuk mempermudah masyarakat menjalani hidup sehat, sementara penelitian tentang aplikasi NutriGrow merancang antarmuka pemantauan gizi balita di Posyandu menggunakan metode \emph{Design Thinking}. Penelitian serupa lainnya merancang aplikasi pendamping diet dengan metode \emph{Design Thinking} dan menguji hasilnya menggunakan \emph{User Experience Questionnaire} (UEQ).

Penelitian-penelitian tersebut menunjukkan bahwa perancangan aplikasi pemantauan asupan gizi telah banyak dilakukan, mulai dari instrumen bagi tenaga ahli gizi, aplikasi pemantauan gizi harian secara umum, hingga aplikasi yang menyasar kelompok pengguna khusus seperti balita di Posyandu, dengan pendekatan desain yang umumnya menggunakan metode \emph{Design Thinking}. Namun demikian, sejauh penelusuran yang dilakukan, penerapan pendekatan \emph{User-Centered Design} (UCD) yang secara sistematis menempatkan konteks penggunaan dan kebutuhan pengguna sebagai titik awal perancangan, khususnya pada konteks pemantauan asupan gizi selama bulan puasa dengan karakteristik waktu sahur yang sangat singkat dan kondisi pengguna yang mudah lelah, belum banyak dieksplorasi secara khusus. Berdasarkan masalah tersebut, dapat disimpulkan bahwa tantangan mahasiswa dalam memenuhi kebutuhan asupan gizi selama bulan puasa tidak hanya berkaitan dengan pengetahuan mengenai gizi seimbang, tetapi juga dengan keterbatasan waktu, perhatian, dan proses pencatatan serta pemantauan asupan pada waktu sahur dan berbuka. Oleh karena itu, tugas akhir ini berfokus pada perancangan desain interaksi pemantauan asupan gizi berbasis \emph{User-Centered Design} (UCD) yang mempertimbangkan kebutuhan, perilaku, dan konteks penggunaan mahasiswa selama bulan puasa, sehingga rancangan yang dihasilkan dapat mendukung proses pencatatan dan pemantauan asupan secara praktis sesuai dengan kebutuhan pengguna.
% --- Rumusan Masalah ---
\section{Rumusan Masalah}
Berdasarkan penjelasan di atas, dapat dirumuskan beberapa permasalahan utama, yaitu rendahnya kebiasaan memantau asupan gizi, keterbatasan fitur aplikasi gizi yang ada dalam konteks bulan puasa, serta kebutuhan pemenuhan prinsip desain interaksi yang sederhana dan tidak membebani pengguna di tengah keterbatasan waktu sahur. Berdasarkan masalah tersebut, dirumuskan beberapa pertanyaan penelitian sebagai berikut:
\begin{enumerate}
    \item Bagaimana konteks penggunaan dan karakteristik kebutuhan pengguna dalam memantau asupan gizi selama bulan puasa?
    \item Bagaimana merancang desain interaksi aplikasi pemantauan asupan gizi yang sesuai dengan kebutuhan pengguna, keterbatasan waktu, serta kondisi fisik pengguna saat sahur dan berbuka?
    \item Bagaimana tingkat \textit{perceived usability} pengguna terhadap desain interaksi aplikasi pemantauan asupan gizi yang dihasilkan berdasarkan evaluasi \textit{heuristik} dan pengujian \textit{usability} menggunakan SUS?
\end{enumerate}

% --- Tujuan ---
\section{Tujuan}
Berdasarkan rumusan masalah yang dijelaskan pada bagian sebelumnya, tujuan akhir dari tugas akhir ini adalah membuat \textit{high-fidelity prototype} aplikasi pemantauan asupan gizi berbasis \textit{mobile} dengan pendekatan \textit{User-Centered Design} untuk membantu mahasiswa memantau asupan gizi selama bulan puasa.
% --- Batasan Masalah ---
\section{Batasan Masalah}
Mengacu pada identifikasi sebelumnya, terdapat beberapa batasan yakni sebagai berikut:
\begin{enumerate}
    \item Luaran tugas akhir ini berupa prototipe interaktif aplikasi pemantauan asupan gizi berbasis \textit{mobile}. Penelitian ini tidak mencakup pengembangan aplikasi hingga siap digunakan.
    \item Analisis kebutuhan pengguna dibatasi pada mahasiswa Muslim di wilayah Bandung yang menjalankan puasa selama bulan puasa. Oleh karena itu, hasil penelitian tidak dimaksudkan untuk mewakili kondisi seluruh mahasiswa di Indonesia.
    \item Evaluasi dibatasi pada \textit{heuristic evaluation} dan \textit{usability testing} menggunakan SUS terhadap prototipe. Penelitian ini tidak menganalisis perubahan perilaku makan, perubahan status gizi, atau efektivitas kesehatan pengguna dalam jangka panjang. Aplikasi yang dirancang tidak ditujukan untuk program penurunan berat badan.
\end{enumerate}

% --- Metodologi Pengerjaan TA ---
\section{Metodologi}

Metodologi tugas akhir ini disusun berdasarkan kerangka \textit{User-Centered Design} (UCD) sebagai bagian dari pendekatan \textit{Human-Centered Design} (HCD) yang ditetapkan dalam ISO 9241-210:2010 (International Organization for Standardization, 2010). Proses HCD terdiri atas lima langkah yang digunakan secara bertahap. Hasil evaluasi pada langkah terakhir dapat menjadi dasar untuk memperbaiki langkah-langkah sebelumnya.

\begin{enumerate}
    \item \textbf{Merencanakan Proses \textit{Human-Centered Design}} \\
    Langkah awal proses HCD berfokus pada perencanaan menyeluruh sebelum tahap penggalian konteks penggunaan dilakukan. Perencanaan ini meliputi penetapan ruang lingkup, tujuan penelitian, target pengguna, serta rencana pelaksanaan hingga tahap evaluasi. Target pengguna dalam penelitian ini adalah mahasiswa Muslim yang menjalankan ibadah puasa. Metode pengumpulan data dan jadwal penelitian disusun sebagai pedoman untuk tahap berikutnya.

    \item \textbf{Memahami Konteks Penggunaan} \\
    Pada langkah kedua, dilakukan penggalian kebutuhan pengguna dan konteks penggunaan dengan menerapkan \textit{user and context analysis}, yang mencakup empat elemen, yaitu pengguna (\textit{users}), tugas dan tujuan (\textit{tasks and goals}), sumber daya (\textit{resources}), serta lingkungan (\textit{environment}) (Geis \& Tesch, 2019, dalam Borja Rivero, 2023). Elemen pengguna mengidentifikasi target penelitian, yaitu mahasiswa Muslim yang menjalankan ibadah puasa di wilayah Bandung, beserta karakteristik dan pengalaman mereka dalam mengatur konsumsi makanan selama berpuasa. Elemen tugas dan tujuan mengidentifikasi aktivitas pengguna, seperti pencatatan menu sahur sebelum memulai aktivitas sehari-hari, baik untuk menyelesaikan aktivitas secara efisien maupun untuk memberikan rasa tenang karena kebutuhan gizi telah direncanakan.

    Elemen sumber daya menentukan alat yang digunakan pengguna untuk mengoperasikan sistem, yaitu \textit{smartphone}, serta keterbatasan sumber daya yang bersifat habis pakai, terutama waktu karena jendela waktu yang sempit, khususnya saat sahur. Elemen lingkungan mencakup kondisi penggunaan pada waktu dan situasi yang berbeda, seperti waktu sahur dan sore hari saat berbuka, serta kondisi fisiologis pengguna seperti rasa lapar, haus, dan kantuk yang dapat memengaruhi kondisi kognitif pengguna. Pemahaman terhadap keempat elemen ini diperoleh melalui wawancara dan survei kepada mahasiswa yang menjalankan ibadah puasa serta didukung studi literatur. Hasilnya menjadi acuan dalam menentukan kebutuhan pengguna pada langkah selanjutnya.

    \item \textbf{Menentukan Kebutuhan Penggunaan} \\
    Berdasarkan pemahaman terhadap pengguna dan konteks penggunaan yang telah diperoleh sebelumnya, langkah ketiga ini merumuskan kebutuhan pengguna (\textit{specify the usage requirements}) dengan bantuan \textit{Hierarchical Task Analysis} (HTA) untuk menguraikan aktivitas pengguna menjadi tugas dan subtugas secara terstruktur. Hasil analisis tugas ini dikombinasikan dengan pemahaman konteks penggunaan untuk merumuskan kebutuhan yang perlu dipenuhi oleh desain, yaitu agar pengguna dapat melakukan aktivitas pengelolaan menu dan pencatatan konsumsi secara efisien tanpa menambah beban waktu dan kognitif, terutama pada waktu sahur dan berbuka.

    \item \textbf{Mengembangkan Solusi Desain} \\
    Kebutuhan pengguna yang telah dirumuskan pada langkah sebelumnya diterjemahkan ke dalam rancangan interaksi secara bertahap pada langkah keempat, mulai dari \textit{low-fidelity prototype} berupa \textit{wireframe} hingga \textit{high-fidelity prototype} dalam bentuk rancangan antarmuka \textit{mobile}. Luaran dari langkah ini berupa prototipe desain interaksi, bukan aplikasi fungsional yang telah diimplementasikan secara penuh, sesuai dengan ruang lingkup penelitian.

    Konten terkait kebutuhan gizi yang ditampilkan dalam desain, seperti informasi kategori dan porsi menu sahur maupun berbuka, mengacu pada pedoman Gizi Seimbang dari Kementerian Kesehatan Republik Indonesia. Pedoman tersebut digunakan sebagai acuan dalam penyusunan konten desain, bukan sebagai hasil analisis atau perhitungan kebutuhan gizi oleh peneliti.

    \item \textbf{Mengevaluasi Solusi Desain} \\
    Langkah terakhir mengevaluasi solusi desain (\textit{evaluate design solutions against user requirements}) untuk mengetahui sejauh mana rancangan yang dikembangkan telah memenuhi kebutuhan pengguna dan tujuan penelitian. Evaluasi dilakukan melalui dua pendekatan, yaitu evaluasi oleh pihak yang memahami bidang UI/UX menggunakan metode seperti \textit{heuristic evaluation} untuk menilai kemudahan, kejelasan, dan konsistensi interaksi, bukan untuk menilai kandungan atau kebutuhan gizi; serta evaluasi langsung kepada kelompok pengguna target, yaitu mahasiswa Muslim yang menjalankan ibadah puasa, melalui \textit{usability testing} dan kuesioner seperti \textit{System Usability Scale} (SUS) dan/atau \textit{User Experience Questionnaire} (UEQ).

    Hasil evaluasi ini menjadi dasar untuk menentukan apakah diperlukan perbaikan pada kebutuhan pengguna maupun solusi desain. Apabila ditemukan ketidaksesuaian atau permasalahan, proses dapat kembali pada tahap penentuan kebutuhan pengguna atau pengembangan solusi desain untuk dilakukan perbaikan.
\end{enumerate}